\documentclass[10pt,letterpaper]{article}
\usepackage{amsmath,amssymb}
\title{Household Budget Worksheet and Commentary}
\author{Rosalind Faircloth}
\date{July 7, 2016}
\begin{document}
\maketitle

\section{Monthly Summary}
A careful budget is the foundation of financial independence. The table
below lists typical categories of expense and the amount allocated to each.

\begin{table}[htbp]
\centering
\begin{tabular}{|l|c|r|}
\hline
\multicolumn{3}{|c|}{\textbf{Monthly Budget}} \\
\hline
Item & Share & Amount \\
\hline
Housing & 35\% & 1400.00 \\
Food & 15\% & 600.00 \\
Transportation & 10\% & 400.00 \\
Savings & 20\% & 800.00 \\
Miscellaneous & 20\% & 800.00 \\
\hline
\end{tabular}
\caption{Proposed allocation for a monthly income of 4000.}\label{tab:budget}
\end{table}

\bigskip
\section{Reasoning}
The shares in Table~\ref{tab:budget} satisfy
$\sum_{i=1}^{5} s_i = 1$, so no income is left unassigned. If savings
earn interest at an annual rate $r$, compounded yearly, then after $t$ years
a monthly deposit $m$ accumulates to roughly
\[
A = 12m \cdot \frac{(1+r)^t - 1}{r}.
\]

\subsection{Unexpected Costs}
Everyone should maintain an emergency fund. Experts differ about the proper
amount, but three to six months of expenses is a frequent recommendation.

\begin{flushleft}
\textbf{Reminder:} review the budget every quarter, and adjust the
miscellaneous category first.
\end{flushleft}

\vspace{2cm}
\subsection*{Signature}
\begin{flushright}
R.~Faircloth\\
Treasurer
\end{flushright}
\end{document}
